% ભાગ: first-order-logic
% પ્રકરણ: models-theories
% વિભાગ: set-theory

\documentclass[../../../include/open-logic-section]{subfiles}

\begin{document}

\olfileid{fol}{mat}{set}

\olsection{ગણોનો સિદ્ધાંત}

લગભગ આખું ગણિત ગણોના સિદ્ધાંતમાં વિકસાવી શકાય છે. આ સિદ્ધાંતમાં
ગણિત વિકસાવવા માટે ઘણી બાબતો જરૂરી છે. પ્રથમ, સંબંધ~$\in$ માટે
સ્વયંસિદ્ધિઓનો ગણ જોઈએ. જુદાં જુદાં સ્વયંસિદ્ધિ-તંત્રો વિકસાવવામાં
આવ્યાં છે અને કેટલીક વાર~$\in$ના તેમના ગુણધર્મો પરસ્પર વિરોધી પણ હોય
છે. પસંદગીના સ્વયંસિદ્ધિ સહિતનું ઝર્મેલો--ફ્રેન્કેલ ગણસિદ્ધાંત,
$\Log{ZFC}$ તરીકે ઓળખાતું સ્વયંસિદ્ધિ-તંત્ર, તેમાં વિશેષ સ્થાન ધરાવે
છે: તે અત્યાર સુધીમાં સૌથી વધુ વપરાતું અને અભ્યાસ કરાતું તંત્ર છે,
કારણ કે તેનાં સ્વયંસિદ્ધિઓ ગણિતજ્ઞો જે લગભગ બધી બાબતો સાબિત કરી શકવાની
અપેક્ષા રાખે છે તે સાબિત કરવા પૂરતાં નીવડે છે. પરંતુ આ સ્થાપિત કરીએ તે
પહેલાં ગણિતજ્ઞો જે બાબતો વ્યક્ત કરવા માગે છે તે આપણે વ્યક્ત પણ કેવી
રીતે કરી શકીએ એ સ્પષ્ટ કરવું જરૂરી છે. શરૂઆતમાં જ મુશ્કેલી દેખાય છે:
ભાષામાં કોઈ !!{constant}s કે !!{function}s નથી. તેથી ચોક્કસ ગણો (જેમ
કે~$\emptyset$ અથવા~$\Nat$), ગણો પરની ક્રિયાઓ (જેમ કે~$X\cup Y$ અને
$\Pow{X}$), અને ગણો સિવાયની વસ્તુઓ સામેલ કરતી સંબંધો અને વિધેયો જેવી
રચનાઓ વિશે કેવી રીતે વાત કરવી એ પ્રથમ નજરે અસ્પષ્ટ લાગે છે.

શરૂઆત કરીએ તો, ``નો ઘટક છે'' એ આપણને રસ હોય એવો એકમાત્ર સંબંધ નથી;
``નો ઉપગણ છે'' લગભગ એટલો જ અગત્યનો લાગે છે. પરંતુ ``નો ઉપગણ છે''ને
``નો ઘટક છે''ના પદોમાં \emph{વ્યાખ્યાયિત} કરી શકીએ છીએ. તેના માટે
ગણસિદ્ધાંતની ભાષામાં એવું !!a{formula}~$!A(x,y)$ શોધવું પડે જે ગણોની
જોડી~$\tuple{X,Y}$થી ત્યારે અને માત્ર ત્યારે સંતોષાય જ્યારે
$X\subseteq Y$. પરંતુ~$X$ના બધા !!{element}s,~$Y$ના પણ !!{element}s
હોય ત્યારે અને માત્ર ત્યારે~$X$,~$Y$નો ઉપગણ છે. તેથી~$\subseteq$ને આ
સૂત્ર વડે વ્યાખ્યાયિત કરી શકીએ:
\[
\lforall[z][(z \in x \lif z \in y)]
\]
હવે જ્યારે પણ કોઈ સૂત્રમાં સંબંધ~$\subseteq$ વાપરવો હોય, ત્યારે તેના
બદલે આ સૂત્ર વાપરી શકીએ (જ્યાં જરૂરી હોય ત્યાં~$x$ અને~$y$ને યોગ્ય
રીતે બદલીને અને બંધ ચલ~$z$નું નામ બદલીને). દાખલા તરીકે, ઘટકો દ્વારા
ગણોનું નિર્ધારણ કહે છે કે કોઈ પણ ગણો~$x$ અને~$y$ એકબીજામાં સમાયેલા
હોય, તો~$x$ અને~$y$ એક જ ગણ હોવા જોઈએ. તેને
$\lforall[x][\lforall[y][((x\subseteq y\land y\subseteq x)\lif x=y)]]$
વડે, અથવા ઉપરની વ્યાખ્યા વડે~$\subseteq$ને બદલીને નીચે પ્રમાણે વ્યક્ત
કરી શકાય:
\[
\lforall[x][\lforall[y][((\lforall[z][(z \in x \lif z \in y)] \land
    \lforall[z][(z \in y \lif z \in x)]) \lif x = y)]].
\]
હકીકતમાં આ~$\Log{ZFC}$નાં સ્વયંસિદ્ધિઓમાંનું એક, ``ઘટકો દ્વારા
નિર્ધારિત સમાનતાનું સ્વયંસિદ્ધિ,'' છે.

$\emptyset$ માટે કોઈ !!{constant} નથી, પરંતુ ``$x$ ખાલી છે''ને
$\lnot\lexists[y][y\in x]$ વડે વ્યક્ત કરી શકીએ છીએ. પછી
``$\emptyset$ અસ્તિત્વ ધરાવે છે'' એ !!{sentence}
$\lexists[x][\lnot\lexists[y][y\in x]]$ બને છે. આ~$\Log{ZFC}$નું બીજું
સ્વયંસિદ્ધિ છે. (નોંધો કે ઘટકો દ્વારા નિર્ધારિત સમાનતાનું સ્વયંસિદ્ધિ
માત્ર એક જ ખાલી ગણ હોઈ શકે એમ ફલિત કરે છે.) ગણસિદ્ધાંતની ભાષામાં
$\emptyset$ વિશે વાત કરવી હોય ત્યારે આપણે ``એવો ખાલી ગણ છે કે~\dots''
એ રીતે લખીએ. દાખલા તરીકે,~$\emptyset$ દરેક ગણનો ઉપગણ છે એ વ્યક્ત કરવા
આ રીતે લખી શકીએ:
\[
\lexists[x][(\lnot\lexists[y][y \in x] \land \lforall[z][x \subseteq
    z])]
\]
જ્યાં અલબત્ત~$x\subseteq z$ને પણ તેની વ્યાખ્યા વડે બદલવું પડે.

$X\cup Y$ અને~$\Pow{X}$ જેવી ગણો પરની ક્રિયાઓની વાત કરવા આ જ પ્રકારની
યુક્તિ વાપરવી પડે છે. ગણસિદ્ધાંતની ભાષામાં વિધેય-સંકેતો નથી, પરંતુ
વિધેયાત્મક સંબંધો~$X\cup Y=Z$ અને~$\Pow{X}=Y$ને નીચે પ્રમાણે વ્યક્ત
કરી શકીએ:
\begin{align*}
& \lforall[u][((u \in x \lor u \in y) \liff u \in z)]\\
& \lforall[u][(u \subseteq x \liff u \in y )]
\end{align*}
કારણ કે~$X\cup Y$ના !!{element}s બરાબર એ ગણો છે જે~$X$ના
!!{element}s અથવા~$Y$ના !!{element}s છે, અને~$\Pow{X}$ના !!{element}s
બરાબર~$X$ના ઉપગણો છે.
છતાં આથી આપણે~$x\cup y$ અથવા~$\Pow{x}$ને પદ તરીકે વાપરી શકતા નથી;
$X\cup Y=Z$ અને~$\Pow{X}=Y$ સંબંધોને વ્યાખ્યાયિત કરતાં આખાં
!!{formula}s જ વાપરી શકીએ છીએ. હકીકતમાં, આ સંબંધો ક્યારેય સંતોષાય છે
એટલે કે યોગગણો અને ઘાતગણો હંમેશાં અસ્તિત્વ ધરાવે છે એવું હજી આપણને
ખબર નથી. દાખલા તરીકે, !!{sentence}
$\lforall[x][\lexists[y][\Pow{x}=y]]$ એ~$\Log{ZFC}$નું બીજું
સ્વયંસિદ્ધિ (ઘાતગણનું સ્વયંસિદ્ધિ) છે.

હવે ક્રમિત જોડીઓ કે વિધેયોની વાતનું શું? અહીં ક્રમિત જોડીઓ અને
વિધેયોને ખાસ પ્રકારના ગણો તરીકે કેવી રીતે વિચારી શકાય તે સમજાવવું પડે.
ક્રમિત જોડી~$\tuple{x,y}$ને ગણ~$\{\{x\},\{x,y\}\}$ તરીકે વ્યાખ્યાયિત
કરવી એ એક રીત છે. પરંતુ પહેલાની જેમ, આ ગણને નામ આપતું કોઈ
!!a{function} દાખલ કરી શકતા નથી; માત્ર સંબંધ~$\tuple{x,y}=z$, એટલે કે
$\{\{x\},\{x,y\}\}=z$, વ્યાખ્યાયિત કરી શકીએ:
\[
\lforall[u][(u \in z \liff (\lforall[v][(v \in u \liff v = x)] \lor
  \lforall[v][(v \in u \liff (v = x \lor v = y))]))]
\]
આ કહે છે કે~$z$ના !!{element}s~$u$ બરાબર એવા ગણો છે જેનો એકમાત્ર
!!{element}~$x$ હોય અથવા જેના એકમાત્ર !!{element}s~$x$ અને~$y$ હોય
(બીજા શબ્દોમાં, એવા ગણો જે~$\{x\}$ કે~$\{x,y\}$ સાથે અભિન્ન હોય).
આ મળ્યા પછી આગળની બાબતો કહી શકીએ, જેમ કે~$X\times Y=Z$:
\[
\lforall[z][(z \in Z \liff \lexists[x][\lexists[y][(x \in
        X \land y \in Y \land \tuple{x, y} = z)]])]
\]

વિધેય~$f\colon X\to Y$ને સંબંધ~$f(x)=y$, એટલે કે જોડીઓના ગણ
$\Setabs{\tuple{x,y}}{f(x)=y}$ તરીકે વિચારી શકાય. પછી ગણ~$f$,~$X$થી
$Y$માં જતું વિધેય છે એમ કહી શકીએ જો (a) તે સંબંધ
$\subseteq X\times Y$ હોય, (b) તે સર્વાગત હોય, એટલે દરેક~$x\in X$
માટે એવું કોઈ~$y\in Y$ હોય કે $\tuple{x,y}\in f$, અને (c) તે
વિધેયાત્મક હોય, એટલે જ્યારે પણ
$\tuple{x,y},\tuple{x,y'}\in f$, ત્યારે $y=y'$ (કારણ કે વિધેયનાં
મૂલ્યો અનન્ય હોવાં જોઈએ). તેથી ``$f$,~$X$થી~$Y$માં જતું વિધેય છે''ને
આ રીતે લખી શકાય:
\begin{align*}
& \lforall[u][(u \in f \lif
   \lexists[x][\lexists[y][(x \in X \land y \in
      Y \land \tuple{x, y} = u)]])] \land {}\\
& \lforall[x][(x \in X \lif
   (\lexists[y][(y \in Y \land \mathrm{maps}(f, x, y))] \land {}\\
&\qquad \lforall[y][\lforall[y'][((\mathrm{maps}(f, x, y) \land
    \mathrm{maps}(f, x, y')) \lif y = y')]]))]
\end{align*}
\sourcecorrection{OLMAT-003}{મૂળની \texttt{set-theory.tex},
પંક્તિઓ~121--125માં વિધેયની વ્યાખ્યાના બંને સાર્વત્રિક વાક્યોમાં
નિષ્કર્ષ શરૂ થાય તે પહેલાં જ પરિમાણકનું દલીલ-કૌંસ બંધ થાય છે. તેથી
સંબંધમાં આવવું અને સર્વાગત-વિધેયાત્મક હોવું જણાવતા નિષ્કર્ષો અનુક્રમે
$u$ અને~$x$ના પરિમાણકની બહાર રહી જાય છે. અહીં બંને નિષ્કર્ષોને તેમના
સાર્વત્રિક પરિમાણકોના વ્યાપમાં પુનઃસ્થાપિત કર્યા છે.}
અહીં~$\mathrm{maps}(f,x,y)$ એ
$\lexists[v][(v\in f\land\tuple{x,y}=v)]$નું સંક્ષિપ્ત સ્વરૂપ છે (આ
!!{formula} ``$f(x)=y$'' વ્યક્ત કરે છે).

હવે $f\colon X\to Y$ !!{injective} છે એ વ્યક્ત કરવું પણ મુશ્કેલ નથી:
\begin{multline*}
 f \colon X \to Y \land \lforall[x][\lforall[x'][((x \in X \land x' \in
     X \land {} \\
 \lexists[y][(\mathrm{maps}(f, x, y) \land \mathrm{maps}(f,
        x', y))]) \lif x = x')]]
\end{multline*}
\sourcecorrection{OLMAT-004}{મૂળની \texttt{set-theory.tex},
પંક્તિઓ~134--136માં એક-એકપણાના સૂત્રના બંને સાર્વત્રિક પરિમાણકો
અસ્તિત્વવાચક પૂર્વધારણા પહેલાં બંધ થાય છે અને અંતે તેમનાં સમાપક કૌંસ
નથી. અહીં આખા શરતી વિધાનને બંને સાર્વત્રિક પરિમાણકોના વ્યાપમાં મૂક્યું
છે.}
વિધેય~$f\colon X\to Y$ !!{injective} ત્યારે અને માત્ર ત્યારે જ્યારે
$f$,~$x,x'\in X$ બંનેને એક જ~$y$ પર મોકલે તો $x=x'$. આ સૂત્રને
$\mathrm{inj}(f,X,Y)$થી સંક્ષિપ્ત કરીએ, તો ગણસિદ્ધાંતની ભાષામાં
કૅન્ટૉરનું પ્રમેય જેટલી બિનતુચ્છ વાત પણ કહી શકીએ:~$\Pow{X}$થી~$X$માં
જતું કોઈ !!{injective} વિધેય નથી:
\[
\lforall[X][\lforall[Y][(\Pow{X} = Y \lif
    \lnot\lexists[f][\mathrm{inj}(f, Y, X)])]]
\]

કોઈને એમ લાગી શકે કે ગણસિદ્ધાંતને દરેક વ્યાખ્યાયિત ગુણધર્મ માટે ગણના
અસ્તિત્વની ખાતરી આપતું બીજું સ્વયંસિદ્ધિ જોઈએ. જો~$!A(x)$ એ મુક્ત
ચલ~$x$ ધરાવતું ગણસિદ્ધાંતનું સૂત્ર હોય, તો !!{sentence}
\[
\lexists[y][\lforall[x][(x \in y \liff !A(x))]].
\]
વિચારી શકીએ. આ !!{sentence} કહે છે કે એવો ગણ~$y$ છે જેના !!{element}s
બરાબર એ બધા~$x$ છે જે~$!A(x)$ને સંતોષે છે. આ પ્રરૂપને
``ગુણધર્મ-ગણરચના સિદ્ધાંત'' કહે છે. તે ખૂબ ઉપયોગી લાગે છે; દુર્ભાગ્યે
તે અસુસંગત છે. $!A(x)\ident\lnot x\in x$ લો; ત્યારે
ગુણધર્મ-ગણરચના સિદ્ધાંત કહે છે
\[
\lexists[y][\lforall[x][(x \in y \liff x \notin x)]],
\]
એટલે કે પોતાના જ !!{element}s ન હોય એવા બધા ગણોના ગણનું અસ્તિત્વ છે.
આવો કોઈ ગણ હોઈ શકે નહિ---આ રસેલનો વિસંવાદ છે. હકીકતમાં,~$\Log{ZFC}$માં
આ સિદ્ધાંતનું મર્યાદિત---અને સુસંગત---સ્વરૂપ, પૃથક્કરણ સિદ્ધાંત, છે:
\[
\lforall[z][\lexists[y][\lforall[x][(x \in y \liff (x \in z \land
      !A(x)))]]].
\]
\sourcecorrection{OLMAT-005}{મૂળની \texttt{set-theory.tex},
પંક્તિઓ~169--170માં પૃથક્કરણ સિદ્ધાંતના અંતે અંદરની સંયોજન અને
$!A(x)$ બંધ થાય છે, પરંતુ દ્વિશરતી વિધાનનું બહારનું ગોળ કૌંસ બંધ થતું
નથી. અહીં ખૂટતું સમાપક કૌંસ ઉમેર્યું છે.}

\begin{prob}
નીચે દર્શાવતી !!a{derivation} આપીને ગુણધર્મ-ગણરચના સિદ્ધાંત અસુસંગત
છે એ બતાવો:
\[
\lexists[y][\lforall[x][(x \in y \liff x \notin x)]] \Proves \lfalse.
\]
પહેલાં $(A\lif\lnot A)\land(\lnot A\lif A)\Proves\lfalse$ બતાવવું
ઉપયોગી થઈ શકે છે.
\end{prob}
\end{document}
